\section{Counting general first-order Taylor--Wiles deformations}

% BEGIN T70B1
\paragraph{T70B1. The two coefficient modules and the real-place factor.}
Take the original arithmetic data of
$\mathrm{FirstOrderDualArithmeticBound}$, including $\mathrm{CohP}$,
the level $V_0.S'=s.S$, the actual $P:\mathrm{TWPl}(r,T)$ and
$\mathrm{LocCFT}(s)$. Retain the nontrivial residual cyclotomic
test, triviality of $x.\rho$ on the indicated pro-$\ell$ kernel,
$x.N4(V_0.\mathrm{frame}.D0)$, the distinct labeled eigenvalues,
and $\mathrm{Detects}(\kappa,\mathrm{range}(T.sg))$. Put
$A=x.K[[t]]$, $A_m=A/t^m$, $Q=|x.K|$, $a=x.a$,
$b=n-a$, $g=\kappa.g$, $q=T.q$ and $d=[L^+:\mathbb Q]$.
All three of $a,b,g$ retain their given values.

Use the coefficient algebra on $1,e_j,e'_j,\eta$ ($1\le j\le g$)
in which $e_i e'_j=\delta_{ij}\eta$ and the remaining products
of two nonconstant basis elements are zero. Write $M$ for
$\mathrm{W1}$ and $M'$ for its paired dual over $A$ or $A_m$.
On the two diagonal blocks $M$ has scalar entries and $M'$ has
$\eta$ entries; on the two off-diagonal blocks $M$ has $g$ entries
in the $e_j$ and $M'$ has $g$ in the $e'_j$. Both have rank
$n+2\nu$ over the coefficient ring, where
\[
       \nu=\binom a2+\binom b2+gab.
\]
Use $R(s)=(1+\sum e_j\kappa_j(s))\rho(s)$ and
$W_v=(1-\sum e_jY_{j,v})w_v$. The cocycle equation for each
$\kappa_j$ and $e_i e_j=0$ give multiplicativity of $R$.
The conjugate-transpose equation follows by the polarized cocycle
condition and the form of $\rho$. At $v$ the infinitesimal
correction to $W_v^{-1}R(s)W_v$ is the sum of the terms
$e_j w_v^{-1}(\kappa_j(s)-\delta Y_{j,v}(s))w_v$ times the
original upper triangular matrix. Thus the local basis makes
$R$ upper triangular. Its diagonal is exactly $\psi_i$, since
these corrections are off block in the ordered local type.
Indeed column $i$ of $w_v$ belongs entirely to the block
specified by $x.\mathrm{ty}(v,i)$, and invertibility gives $a$
columns of the first type and $b$ of the second. Conjugation
by the off-block $e_j$ corrections preserves the pattern of
scalar entries on equal types and $e_j$ entries on unequal
types: mixed block products carry one $e_j$, and two such
factors vanish. The matrices $W_v$ and $W_v^{-1}$ therefore
identify $M$ with that same pattern in the local ordering.

The action on $M$ is $X\mapsto RXR^{-1}$ on the sign kernel and
$X\mapsto-X^t$ for the marked conjugation $c$. On $M'$ it is the
same action multiplied by the residual cyclotomic character
$\chi$. The form identity ensures the semidirect-product relation,
and direct multiplication verifies that both modules are preserved.
In particular a product of an $e_j$ entry with an $e'_j$ entry
lands in the $\eta$ block. At $m>0$ the pairing is
\[
 (X,Y)\longmapsto
      \lambda\bigl([t^{m-1}]\,[\eta]\operatorname{tr}(XY)\bigr),
\]
where $\lambda:x.K\to\mathbb F_\ell$ is a nonzero linear
functional. A single transposed matrix unit with complementary
coefficient $1,\eta$ or $e_j,e'_j$ detects each individual
coordinate by the first-nonzero-coefficient argument of CFO10.
It is a perfect Tate pairing by trace cyclicity and transpose,
including the factor $\chi$ in the dual action. The zero modules
at $m=0$ also pair perfectly.

The arithmetic presentation and the enumeration of its real and
finite places in CFO11--14 use only $s,V_0,T,P$. Their constructions
use the original level isomorphisms, splitting and local maps
of $P$ and apply to every such $P$ here. For a real conjugation
$c_i$ let $S=R(p_N(c_ic_N))$. Its argument has cyclotomic
character one and is sent to its inverse by marked conjugation.
The form identity consequently makes $S$ symmetric invertible.
$S,S^{-1}$ have scalar diagonal blocks and off-block $e_j$
coordinates and multiply $M$ into itself. The action of $c_i$
on $M$ is $X\mapsto-SX^tS^{-1}$. Its fixed elements are $SZ$
for skew $Z\in M$, in bijection with their $\nu$ independent
upper coordinates since $\ell$ is odd. Thus
$|H^0(c_i,M_m)|=Q^{m\nu}$, including when a block is empty.
% END T70B1

% BEGIN T70B2
\paragraph{T70B2. Three uniform invariant bounds and the flagged local count.}
The DVR kernel estimate of CFO8 applies to these finite free
coordinate modules: a common kernel of dimension $r$ over
$\mathrm{Frac}(A)$ has reduction kernels of size at most
$Q^{mr+C}$ for all $m$. Indeed finitely many equations already
define that fraction-field kernel, and the adjugate of a
nonzero rank-size minor bounds the remaining torsion. We
apply it three times.

For the global primal invariants a skew $Z\in M$ commuting
with $R$ has scalar coordinate $Z_0$ commuting with $\rho$.
The spanning condition on $x.\rho(G')$ gives, over the fraction
field, all endomorphisms of the two blocks separately. Their
commutant consists of scalars on each nonempty block and zero
off block, by the block projectors and the matrix units.
Skewness and odd characteristic force $Z_0=0$. Each $e_j$
coordinate now commutes with $\rho$ by the product formula
and vanishes by the same argument. The resulting global bound
is $Q^{C_1}$. The level-to-base group map is onto, including
on sign kernels, so the bound holds at every level.

At the distinguished old subgroup $D_0$, the $e'_j$ coordinates
of a dual invariant satisfy $\chi(s)\rho(s)Z_j=Z_j\rho(s)$
for all $s\in D_0$. They vanish by the stated $N4(D_0)$,
which passes to the fraction field by clearing denominators.
Their vanishing leaves the identical equation on the $\eta$
coordinate. Hence the old local dual invariants have size
at most $Q^{C_0}$. At a listed $s_j=T.sg(j)$, put $P_j=\rho(s_j)$.
It has $n$ distinct eigenvalues in the fraction field. Its
centralizer is commutative of dimension $n$ by diagonalization,
and contains both block projectors. An off-block $e'_h$ coordinate
commuting with $P_j$ therefore commutes with each projector,
forcing it to vanish. After this, the $\eta$ coordinate belongs
to the ordinary $n$-dimensional centralizer. The modular
centralizer in $M'_m$ of $R(s_j)$ has size at most
$Q^{mn+C_j}$. The constants in all three applications precede
the choice of $m$.

Write $d'_v$ for the local ranks in $\mathrm{LocCFT}$ and $d_v$
for the local degrees in the arithmetic presentation. In the
$W_v$ basis an upper triangular allowed $M_m$ has $n$ scalar
diagonal coordinates and $\nu$ strictly upper coordinates:
one for equal types, $g$ for unequal types. All orderings have
the same total. For its cocycles with diagonal zero on $\mu_v$
take the same filtration by increasing distance $j-i$ as in
CFO16. Once smaller distances vanish, conjugation by the two
upper triangular operators has only the summand with row $i$,
column $j$ left in that coordinate, with scalar multiplier
$\psi_i/\psi_j$. Both diagonal entries of the operators are
scalar, so this holds separately for all $g$ coordinates at an
unequal-type entry. The diagonal factors are homomorphisms
vanishing on $\mu_v$. Each has at most $Q^{m(d'_v+1)}$
possibilities by the local abelianization in CFO15. For a strict
upper entry $x.N1$ gives $\psi_i/\psi_j\ne\chi$. Scalar local
duality and Euler characteristic in CFO15 bound its cocycles by
$Q^{m(1+d_v)+c_{ij}}$. In particular the whole group $Z_{b,v}$
is finite, with
\[
 |Z_{b,v}|\le
 Q^{m\{n(d'_v+1)+\nu(1+d_v)\}+C_v}.
\]
All constants $C_v$ include the indicated multiplicities.
The upper triangular submodule has size $Q^{m(n+\nu)}$.
Its coboundaries lie in $Z_{b,v}$, since upper triangular
conjugation fixes its diagonal. Dividing by the invariant
kernel on this submodule, bounded by the full invariant
kernel, gives
$|B_v^1+Z_{b,v}|/|M_m|\le |Z_{b,v}|/Q^{m(n+\nu)}$ exactly
as in CFO17.
% END T70B2

% BEGIN T70B3
\paragraph{T70B3. Uniform control of the actual dual Selmer classes.}
Let $\Gamma_0$ denote the base group, $G_0$ its sign kernel
and $p_N:\Gamma_N\to\Gamma_0$ the level map. For $s\in G_0$
write $\beta(s)Z=\chi(s)R(s)ZR(s)^{-1}$ on $M'$, and
$\beta(c)Z=-\chi(c)Z^t$. Use the linear functional
\[
        L_j(z)=[\eta]\operatorname{tr}\bigl(z(s_j)R(s_j)\bigr).
\]
In particular it includes the cross terms between the $e'_h$
of $z$ and $e_h$ of $R$.

Over $A$, take the module of functions $G_0\to M'$ continuous
modulo every $t^u$, and the map recording the defects of the
twisted cocycle, the conjugation identity $z(csc^{-1})=\beta(c)z(s)$,
and all $L_j$. This is a linear map into a product of separated
torsion-free modules. Its source is complete, separated and
torsion-free by coordinatewise limits in the finite free module
$M'$; division by $t^u$ preserves continuity, and every locally
constant $M'_u$-valued function has a lift of this kind. The
mod-$t$ kernel of the equations is finite by the finite-cocycle
argument of CFO2: on the open kernel of the finite linear action
the coordinates of a cocycle are continuous homomorphisms to
$\mathbb F_\ell$, finite in number by $\Phi_\ell(\Gamma_0)$.
Finitely many cosets complete the count. Here $\Phi_\ell(\Gamma_0)$
is supplied by $\mathrm{CohP}$ and $V_0.S'=s.S$ exactly as in
CAP1--2. The normalized coboundaries
$B=\{\delta Y:\beta(c)Y=Y\}$ lie in the kernel; $\chi(s_j)=1$
and trace cyclicity give $L_j(\delta Y)=0$.
The original $\mathrm{Detects}$ condition says that a power of
$t$ takes each integral element of this kernel into $B$. The
complete-DVR approximation lemma CFO3 now gives one $C_2$
with
\[
 \#\{z\in\mathrm{TwZ}(A_m): L_j(z)=0\ \text{for all }j\}
       \le Q^{C_2}\#\{\delta Y\text{ on }G_0:
                        Y\in M'_m,\ \beta(c)Y=Y\}
\]
for all $m$. The reduction of an integral coboundary is contained
in the group on the right. This argument uses the actual
$\eta$-trace in the assumed detection condition.

It remains to connect these equations to the added local
conditions. Write $H=T.H$. By its definition as an intersection
of open normal subgroups and compactness (CFO5),
$\mathrm{PL}_\ell(H)\subset\ker(\rho\bmod t^m)$ implies that
every $h\in H$ has some $\ell$-power whose image under $\rho$
is one modulo $t^m$. Its image under $R$ then has the form
$1+U$, $U$ in the square-zero $e_j$ coordinates. Hence
$(1+U)^\ell=1$. Moreover $\chi|_H=1$. All the linear actions
in question therefore have $\ell$-power order on $H$ at
level $m$. The finite affine action belonging to any modular
cocycle has this same property there, since an extra $\ell$th
power kills a translation. Its open kernel, or the intersection
of its finitely many conjugates when starting on $G_0$, contains
$\mathrm{PL}_\ell(H)$: normality of $H$ ensures the $\ell$-power
property in every conjugate. This works for the finite matrix
representation $R\bmod t^m$ itself as well.

Fix $m$. Apply $T.\mathrm{ph\_lim}$ to these normal kernels on
the base group, also including $\ker\chi$. There are only
finitely many modular twisted cocycles by the argument above,
and finitely many $s_j$. Thus eventually in $N$ the full matrices
$R(p_N(\mathrm{ph}_{Nj}))$ and $R(s_j)$ agree, $\chi$ has
value one there, and each modular twisted cocycle has equal
values at these two elements. These choices of $N$ precede
every subsequent choice of cocycle at level $N$.

For a dual Selmer class, its local restriction at an added
place is zero, since the primal condition there is unrestricted.
Normalize a representing cocycle to vanish at $c_N$ by subtracting
$\delta(-z(c_N)/2)$. The same local coboundary at the place
then evaluates its normalized representative at both
$\mathrm{ph}_{Nj}$ and $\mathrm{ta}_{Nj}$. The latter value
is zero since $p_N(\mathrm{ta}_{Nj})=1$. The affine kernel
also contains $\mathrm{PL}_\ell(p_N^{-1}H)$ by the $\ell$-power
argument. The kernel bound $T.\mathrm{ker\_le}$ puts $\ker p_N$
in the affine kernel, so the cocycle descends continuously by
the compact quotient map. Its restriction to $G_0$ satisfies
the twisted cocycle and conjugation equations. The equalities
at Frobenius above replace its value there by a coboundary at
$s_j$, which kills $L_j$ by trace cyclicity. Normalization at
$c_N$ makes restriction to $G_0$ after descent injective on
these representatives. More precisely it is injective on their
classes modulo the normalized coboundaries: any $\delta Y$ on
$G_0$ with $\beta(c)Y=Y$ inflates to the corresponding normalized
level coboundary, and the two sign cosets determine its values.
We have proved that the dual Selmer group has at most $Q^{C_2}$
classes, eventually in $N$ for each $m$.
% END T70B3

% BEGIN T70B4
\paragraph{T70B4. The first-order arithmetic bound for every original point and basis.}
At level $N$ put the local conditions $B_v^1+Z_{b,v}$ above
$\ell$, the cocycles vanishing on prescribed inertia at each
old tame place, and unrestricted classes at the distinguished
place and at the $q$ added places. On prescribed tame inertia
the action of $R$ is trivial: $\rho=1$ there by $x.\mathrm{ine}$
and each $\kappa_j$ is zero there by membership in the given
odd Selmer space. The unramified calculation in CFO11 and CFO14
therefore gives a factor one at each such place in the product
of local cocycle sizes divided by $|M_m|$. At any place
$w\nmid\ell$ with unrestricted classes, the local Euler
identity from $\mathrm{CohP}$ gives
$|Z^1(D_w,M_m)|/|M_m|=|H^0(D_w,M'_m)|$ by the pairing of T70B1.
At the old distinguished place its image is exactly $D_0$.
At an added place it contains $\mathrm{ph}_{Nj}$, which agrees
with $s_j$ in the full finite action on the eventual set of
T70B3. The factors are thus at most $Q^{C_0}$ and
$Q^{mn+C_j}$ by T70B2.

Let $\mathcal S$ denote the resulting primal Selmer group.
The map from cocycles normalized to zero at $c_N$ to functions
on the sign kernel, with inverse $z(sc_N^e)=z(s)$, identifies
$\mathcal S$ with $\mathrm{Sel1}_N/\mathrm{Bd1}_N$: the action
of $c_N$ gives the required negative transpose, the $W_v$ bases
and local coboundaries give exactly the original flag
conditions, and the normalized global coboundaries are given
by skew matrices. This is the argument of CFO18 on the modules
in T70B1. The finite local and global cohomology in $\mathrm{CohP}$
and finiteness of the modules give finite groups of cocycles.

Apply the Greenberg--Wiles identity of the assumed $\mathrm{CohP}$
to this finite Tate pair and these very local conditions,
as in CFO13 and CFO19. In cocycle notation its upper bound is
\[
 |\mathcal S|\prod_{i=1}^d |M_m^{c_i}|
 \le |\mathcal S^\perp|\,|M_m^{\Gamma_N}|
               \prod_{w}\frac{|\mathcal L_w|}{|M_m|},
\]
where $\mathcal L_w$ is the group of local cocycles specified
above. This follows by canceling local invariant orders in
$|\mathcal L_w/B^1_w|=|\mathcal L_w|\,|M_m^{D_w}|/|M_m|$
and using that the global dual invariant order is at least one.
For ordinary places the product factor is bounded in T70B2.
The real-place product is $Q^{m\nu d}$. Since
$\sum_v d'_v=\sum_v d_v=d$, the sum of the ordinary exponents
after dividing by the triangular orders is
\[
 \sum_{v\mid\ell}\!
  \bigl(m\{n(d'_v+1)+\nu(1+d_v)\}+C_v-m(n+\nu)\bigr)
       =m(nd+\nu d)+\sum_v C_v .
\]
Cancel the real-place exponent. The $q$ added terms contribute
$mnq+\sum_j C_j$. The global primal invariants and dual Selmer
classes have the bounds of T70B2 and T70B3. Setting
\[
 C=C_0+C_1+C_2+\sum_v C_v+\sum_j C_j
\]
before $m,N$ are chosen gives
\[
 |\mathrm{Sel1}_N(A_m,\kappa)|
 \le |\mathrm{Bd1}_N(A_m,\kappa)|\,Q^{m(nd+nq)+C}.
\]
For each $m$ this holds eventually in $N$; for $m=0$ both sets
are singletons. This proves the complete
$\mathrm{FirstOrderDualArithmeticBound}$ with its original
quantifier order, $D_0=V_0.\mathrm{frame}.D0$, and unrestricted
$a$ and $\kappa.g$.
% END T70B4
